% Shared preamble for methods.tex (the technical document) and explained.tex
% (its lay companion). Engine: pdflatex (pinned in latexmkrc).
\usepackage[utf8]{inputenc}
\usepackage[T1]{fontenc}
% The kernel's own labels reach the generated fragments verbatim (a
% product recipe is "2f1−f2" with a true minus sign, a pitch is "C♯4"), so
% the characters the shipped vocabulary uses are declared once here for
% both documents rather than rewritten by a generator.
\DeclareUnicodeCharacter{2212}{\ensuremath{-}}
\DeclareUnicodeCharacter{266F}{\ensuremath{\sharp}}
\DeclareUnicodeCharacter{266D}{\ensuremath{\flat}}
\DeclareUnicodeCharacter{2026}{\dots}
\DeclareUnicodeCharacter{2248}{\ensuremath{\approx}}
\DeclareUnicodeCharacter{2264}{\ensuremath{\le}}
\DeclareUnicodeCharacter{2265}{\ensuremath{\ge}}
\DeclareUnicodeCharacter{00D7}{\ensuremath{\times}}
\DeclareUnicodeCharacter{03B4}{\ensuremath{\delta}}
\usepackage{lmodern}
\usepackage{amsmath,amssymb}
\usepackage{booktabs,longtable,array}
\usepackage{graphicx}
\usepackage[margin=24mm]{geometry}
\usepackage{xcolor}
\usepackage[numbers,sort&compress]{natbib}
% hyperref: every cross-reference, citation, glossary term, TOC entry and
% URL in the PDF is a clickable link; bookmarks mirror the sections.
\usepackage[colorlinks=true,linkcolor=blue!50!black,citecolor=blue!50!black,
  urlcolor=blue!50!black,linktoc=all,bookmarksnumbered=true,bookmarksopen=true,
  pdfusetitle,pdfpagemode=UseOutlines,breaklinks=true]{hyperref}
\hypersetup{pdfsubject={PedalScope measurement methods},pdfkeywords={swept sine, harmonic distortion, parity, reproducibility}}
% DOIs in the bibliography resolve as links (plainnat emits \doi{...}).
\providecommand{\doi}{}\renewcommand{\doi}[1]{\href{https://doi.org/#1}{doi:#1}}
\newcommand{\repo}{\url{https://github.com/PhaseDog/PedalScope}}
% The licence line both documents print in their front matter (ruled
% 2026-09-08): the holder as the app's About view states it, CC BY 4.0 for
% the document, MIT for the scripts, the mark excluded. One definition, two
% uses; the two URLs are hyperref links (the pdf job's census counts them).
\newcommand{\licenceline}{%
  \begin{center}\small
  \copyright{} 2026 Richard Hoge. This document is licensed under CC~BY~4.0
  (\url{https://creativecommons.org/licenses/by/4.0/}). The accompanying
  scripts are licensed under MIT
  (\url{https://opensource.org/license/mit}). PedalScope and its mark are
  not part of either grant.
  \end{center}}
\graphicspath{{../figures/}{../../UserGuide/assets/}}
\setlength{\parskip}{4pt plus 1pt}
\setlength{\parindent}{0pt}
% Terms: \term{key} prints the glossary name and links to the entry;
% \termas{key}{text} links custom text. Both are what the terms guard scans.
% GENERATED by gen_docs.py from terms.yaml — do not edit
\expandafter\newcommand\csname term@manifest\endcsname{methods manifest}
\expandafter\newcommand\csname term@parity\endcsname{parity test}
\expandafter\newcommand\csname term@sweep\endcsname{synchronized swept sine}
\expandafter\newcommand\csname term@rate-constant\endcsname{rate constant}
\expandafter\newcommand\csname term@synchronization\endcsname{synchronization condition}
\expandafter\newcommand\csname term@instantaneous-frequency\endcsname{instantaneous frequency}
\expandafter\newcommand\csname term@excited-band\endcsname{excited band}
\expandafter\newcommand\csname term@fundamental\endcsname{fundamental}
\expandafter\newcommand\csname term@harmonic-order\endcsname{harmonic order}
\expandafter\newcommand\csname term@clipping\endcsname{clipping}
\expandafter\newcommand\csname term@harmonic-response\endcsname{harmonic response}
\expandafter\newcommand\csname term@harmonic-pulse\endcsname{harmonic pulse}
\expandafter\newcommand\csname term@deconvolution\endcsname{deconvolution}
\expandafter\newcommand\csname term@inverse-filter\endcsname{inverse filter}
\expandafter\newcommand\csname term@impulse-response\endcsname{impulse response}
\expandafter\newcommand\csname term@extraction-window\endcsname{extraction window}
\expandafter\newcommand\csname term@input-normalized\endcsname{input-normalized}
\expandafter\newcommand\csname term@validity-bound\endcsname{validity bound}
\expandafter\newcommand\csname term@alias-image\endcsname{alias image}
\expandafter\newcommand\csname term@thd\endcsname{total harmonic distortion}
\expandafter\newcommand\csname term@floor\endcsname{measurement floor}
\expandafter\newcommand\csname term@fundamental-presence\endcsname{fundamental presence}
\expandafter\newcommand\csname term@shaping\endcsname{harmonic shaping}
\expandafter\newcommand\csname term@coherence\endcsname{coherence}
\expandafter\newcommand\csname term@scatter\endcsname{read scatter}
\expandafter\newcommand\csname term@averaging\endcsname{sweep averaging}
\expandafter\newcommand\csname term@derotation\endcsname{derotation}
\expandafter\newcommand\csname term@tanh\endcsname{tanh}
\expandafter\newcommand\csname term@phantom\endcsname{phantom capture}
\expandafter\newcommand\csname term@decibel\endcsname{decibel}
\expandafter\newcommand\csname term@dbfs\endcsname{dBFS}
\expandafter\newcommand\csname term@dbc\endcsname{dBc}
\expandafter\newcommand\csname term@transfer-curve\endcsname{transfer curve}
\expandafter\newcommand\csname term@averaged-cycle\endcsname{averaged cycle}
\expandafter\newcommand\csname term@static-curve\endcsname{static curve}
\expandafter\newcommand\csname term@box-normalization\endcsname{box normalization}
\expandafter\newcommand\csname term@lobe-area\endcsname{unsigned lobe area}
\expandafter\newcommand\csname term@fundamental-ellipse\endcsname{fundamental ellipse}
\expandafter\newcommand\csname term@nonlinear-loop-area\endcsname{nonlinear loop area}
\expandafter\newcommand\csname term@phase-removal\endcsname{phase removal}
\expandafter\newcommand\csname term@frame-fit\endcsname{frame fit}
\expandafter\newcommand\csname term@lattice\endcsname{memoryless lattice}
\expandafter\newcommand\csname term@pairing-residual\endcsname{pairing residual}
\expandafter\newcommand\csname term@polarity\endcsname{polarity}
\expandafter\newcommand\csname term@orientation\endcsname{orientation resolution}
\expandafter\newcommand\csname term@coupling-capacitor\endcsname{coupling capacitor}
\expandafter\newcommand\csname term@alignment-error\endcsname{residual alignment error}
\expandafter\newcommand\csname term@memory-verdict\endcsname{memory verdict}
\expandafter\newcommand\csname term@two-tone\endcsname{two-tone measurement}
\expandafter\newcommand\csname term@intermodulation-product\endcsname{intermodulation product}
\expandafter\newcommand\csname term@product-order\endcsname{product order}
\expandafter\newcommand\csname term@analysis-lattice\endcsname{analysis lattice}
\expandafter\newcommand\csname term@reduced-ratio\endcsname{reduced tone ratio}
\expandafter\newcommand\csname term@lattice-spacing\endcsname{lattice spacing}
\expandafter\newcommand\csname term@lattice-snap\endcsname{lattice snap}
\expandafter\newcommand\csname term@coherent-sampling\endcsname{coherent sampling}
\expandafter\newcommand\csname term@analysis-window\endcsname{analysis window}
\expandafter\newcommand\csname term@local-floor\endcsname{local floor}
\expandafter\newcommand\csname term@capture-floor\endcsname{capture floor}
\expandafter\newcommand\csname term@gate\endcsname{presence gate}
\expandafter\newcommand\csname term@product-reading\endcsname{product reading}
\expandafter\newcommand\csname term@coherence-check\endcsname{coherence check}
\expandafter\newcommand\csname term@coherence-detector\endcsname{coherence detector}
\expandafter\newcommand\csname term@pair-statistic\endcsname{sideband pair statistic}
\expandafter\newcommand\csname term@harmonic-role\endcsname{harmonic role}
\expandafter\newcommand\csname term@just-ratio\endcsname{just ratio}
\expandafter\newcommand\csname term@gain-map\endcsname{gain map}
\expandafter\newcommand\csname term@drive-level\endcsname{drive level}
\expandafter\newcommand\csname term@level-lattice\endcsname{level lattice}
\expandafter\newcommand\csname term@drive-window\endcsname{drive window}
\expandafter\newcommand\csname term@map-cell\endcsname{cell}
\expandafter\newcommand\csname term@export-grid\endcsname{export grid}
\expandafter\newcommand\csname term@truncated-column\endcsname{truncated column}
\expandafter\newcommand\csname term@edge-mask\endcsname{edge mask}
\expandafter\newcommand\csname term@delivered-fundamental\endcsname{delivered fundamental}
\expandafter\newcommand\csname term@floor-stamp\endcsname{calibration stamp}
\expandafter\newcommand\csname term@median-read\endcsname{median read}
\expandafter\newcommand\csname term@drive-term\endcsname{drive term}
\expandafter\newcommand\csname term@operative-floor\endcsname{operative floor}
\expandafter\newcommand\csname term@analyzer-residue\endcsname{analyzer residue}
\expandafter\newcommand\csname term@composed-margin\endcsname{composed margin}
\expandafter\newcommand\csname term@corroborated-peak\endcsname{corroborated peak}
\expandafter\newcommand\csname term@peak-tile\endcsname{peak tile}
\expandafter\newcommand\csname term@cleanup-reading\endcsname{cleanup reading}
\expandafter\newcommand\csname term@audition-model\endcsname{audition model}
\expandafter\newcommand\csname term@leakage-map\endcsname{leakage map}
\expandafter\newcommand\csname term@compression-curve\endcsname{compression curve}
\expandafter\newcommand\csname term@probe-grid\endcsname{probe grid}
\expandafter\newcommand\csname term@fine-window\endcsname{fine window}
\expandafter\newcommand\csname term@stepped-tone\endcsname{stepped tone}
\expandafter\newcommand\csname term@settle-window\endcsname{settle region and measurement window}
\expandafter\newcommand\csname term@cycle-snap\endcsname{cycle snap}
\expandafter\newcommand\csname term@one-bin-estimator\endcsname{one-bin estimator}
\expandafter\newcommand\csname term@noise-stamp\endcsname{noise stamp}
\expandafter\newcommand\csname term@noise-gate\endcsname{per-step noise gate}
\expandafter\newcommand\csname term@incremental-slope\endcsname{incremental slope}
\expandafter\newcommand\csname term@hinge-fit\endcsname{hinge fit}
\expandafter\newcommand\csname term@knee\endcsname{knee}
\expandafter\newcommand\csname term@bound-zone\endcsname{floor-bound zone}
\expandafter\newcommand\csname term@resolution-band\endcsname{resolution band}
\expandafter\newcommand\csname term@unity-guard\endcsname{unity guard}
\expandafter\newcommand\csname term@noise-bound\endcsname{noise bound}
\expandafter\newcommand\csname term@floor-class\endcsname{floor class}
\expandafter\newcommand\csname term@floor-runs\endcsname{floor runs}
\expandafter\newcommand\csname term@null-run-reference\endcsname{null-run reference}
\expandafter\newcommand\csname term@floor-source\endcsname{floor source}
\expandafter\newcommand\csname term@summed-presence\endcsname{summed presence reading}

% The machine floor (ruled 2026-09-22, py/machine_floor.py): a measured
% deviation under \machineFloor{} is printed as the bound "≤ floor", and
% the sentence each parity section states once is defined here, once.
% GENERATED by gen_docs.py from machine_floor.py — do not edit
\newcommand{\machineFloor}{1e-10}
\newcommand{\machineFloorWords}{a ten-billionth}
\newcommand{\sidecarFigures}{12}

\newcommand{\machineFloorNote}{A measured figure under \machineFloor{} is
printed as the bound $\le$\,\machineFloor: below it the digits are the
machine's last-bit arithmetic rather than the method's, and two machines
running the same code do not share them. A stated tolerance keeps its
digits.}
\newcommand{\term}[1]{\hyperlink{term:#1}{\csname term@#1\endcsname}}
\newcommand{\termas}[2]{\hyperlink{term:#1}{#2}}
\newcommand{\code}[1]{\texttt{#1}}
\newcommand{\dB}{\ensuremath{\,\mathrm{dB}}}
